\label{sec:discussion}

\section{Summary of Findings}

Adaptive per-round privacy budget allocation cannot outperform uniform
allocation in clipped DP-FL, because clipping removes the very signal variation
the approach is designed to exploit. We prove this
(Theorem~\ref{thm:uniform}), verify it across 13 allocations from three
unrelated families ordered by a single closed-form scalar, and show that the
bound applies to published learned allocators because their state space confines
them to fixed schedules. Training the agent and running it confirms the
prediction: it learns a non-uniform schedule and loses to uniform by $-0.043$
mAP over paired held-out trials. The same inequality closes the analogous
problems of allocating a budget across layers and across clients, neither of
which requires a separate experiment.

\vspace{10pt}

On the security side, the same high dimensionality that makes DP costly for
detection also makes it protective. Update norms concentrate as $1/\sqrt{2D}$,
attacker amplitude is capped by the resulting narrow envelope, and the observed
degradation under adaptive poisoning at $30\%$ malicious clients is $2.7\%$
relative. Detection-based defences neutralise what remains; classical robust
aggregators cost more than the attacks do.

\section{Implications for Practice}

Three recommendations follow for practitioners building differentially private
federated systems.

\begin{enumerate}[label=\arabic*.]
  \item \textbf{Spend the privacy budget uniformly} and invest engineering
  effort elsewhere. Uniform is not merely a reasonable default but the analytic
  optimum under per-sample clipping, and the effort required to search for a
  better schedule is better spent on the architecture or the data partition.
  \item \textbf{Do not use adaptive optimisers with DP-SGD.} Adam normalises the
  privacy parameter away: a sixteen-fold change in $\sigma$ moved accuracy by
  $0.016$, while the privacy accountant continued to report the intended
  $\varepsilon$. This is a silent failure that inflates apparent utility at a
  claimed privacy level.
  \item \textbf{Prefer detection over robust aggregation when DP is present},
  and report the false-positive cost. Averaging is what suppresses DP noise, so
  replacing it with a robust statistic forfeits that suppression; in our setting
  coordinate-wise median cost an order of magnitude more than any attack tested.
\end{enumerate}

\section{Implications for the Literature}

The results suggest that reported gains from adaptive allocation, often around
$0.3\%$ relative and published without error bars, are within the noise of the
experimental designs used to obtain them. The CRN protocol used here reaches at
$n=20$ the precision that $n\approx150$ would give without it, and at that
precision the effect is absent.

\vspace{10pt}

We are careful about what this does and does not say. The claim is not that the
original work is mistaken about what it observed. The claim is that the
observation is not separable from noise at the precision used to obtain it, and
that a mechanism argument, Theorem~\ref{thm:uniform} together with the
measurement in Table~\ref{tab:clipbinds} showing that clipping binds by
$8\times$ even in the source configuration, predicts that no separable effect
should exist.

\vspace{10pt}

A negative result of this kind is only as strong as the effort spent trying to
make the positive result appear. That effort is documented rather than asserted:
Chapter~5 reports a pre-registered search over three schedule families, a mild
front-load chosen specifically to be cheap in variance, early stopping, and a
ladder of privacy budgets down to the collapse threshold, with all outcomes
reported including those that initially looked favourable. Four separate
single-run results that appeared positive did not survive replication, and each
is reported.

\section{Limitations}
\label{sec:limitations}

Results are obtained on a clear-sky subset partitioned by recording session, and
absolute numbers do not transfer to cluttered scenes, where the detector sits
near the floor. We note that this choice favours the alternative hypothesis: it
is the condition under which a non-uniform schedule has the most room to show an
advantage.

\vspace{10pt}

Theorem~\ref{thm:uniform} assumes that clipping binds, verified here at
$27\times$ for the detector and $8\times$ for the source configuration, and that
utility is monotone in accumulated variance. The latter assumption measurably
weakens near the collapse threshold, where the penalty for deviating from
uniform shrinks with slope $-0.021$ to $-0.029$ against $\log\rho$, though never
enough to change the conclusion. Configurations in which clipping does not bind,
for instance a very large $C$ with small gradients, fall outside the theorem;
there the premise would be live, though such configurations also forfeit much of
the privacy guarantee that motivates clipping.

\vspace{10pt}

The poisoning study covers two attacks at $m\le30\%$ malicious clients, at a
single noise level, on the clear-sky subset. The cross-client extension of the
theorem is proved but not verified empirically. The implementation port to
CIFAR-10 described in Section~\ref{sec:provenance} is single-run per arm and no
conclusion rests on it.

\section{Future Work}
\label{sec:future}

The one policy class our bound does not cover is an agent whose state includes
model feedback, such as the current loss or realised gradient norm, and which
can therefore react to the noise actually drawn instead of committing to a
schedule in advance. Theorem~\ref{thm:uniform} constrains the expected
accumulated variance of any fixed schedule, but a controller that observes the
realised trajectory is optimising a different, larger object. Whether
exploitable per-run variance exists at all is unmeasured, and is the natural
next question; the online noise-modulation family exemplified by DNO-RL
\cite{WangDNORL} is the closest existing point of comparison.

\vspace{10pt}

On the security side, an attacker that optimises direction against the accuracy
detector, rather than simply reversing it, would be a stronger test than the one
used here, since the accuracy detector is the component that carries the mix
defence. Extending the envelope analysis across model scales would also
establish where the camouflage premise crosses from live to inert, which this
work locates only at its two endpoints.

\section{Conclusion}

The contribution is a negative result with a proof and a positive result with a
mechanism. Uniform privacy-budget allocation is optimal, not merely a reasonable
default, and the learned allocators proposed to improve on it are bounded by
that optimum because their state space confines them to the class of fixed
schedules. The same dimensional scaling that makes differential privacy
expensive for detection also makes the resulting system difficult to poison, by
concentrating the honest update norm and capping the amplitude available to an
attacker. Both findings were obtained in a regime constructed specifically to
give the alternatives every chance to succeed.
